\section{Local comparison}
\label{app:local}

This section proves the local implication used after
Sections~\ref{app:candidate} and \ref{app:global}.
All four profiles use the intrinsic residual and physical normal estimates.
For the odd profiles the complete unsigned comparison suffices.
For $C=3$ we use the stress-coordinate mean estimate with the complete
unsigned remainder. For $C=1$ we use the refined row coefficients and
the minimum of the complete signed alternatives and the complete unsigned
bound.
All derivatives are real derivatives on complex vector spaces, and
$\langle H,J\rangle=\Rea\tr(H^*J)$. Matrix norms without subscripts are
operator norms. Frame distances are measured in the Frobenius metric on
the original variables.

\subsection{Inputs, gauge, and the normal decomposition}
\label{app:local:inputs}

Let $a(i)\in\{1,2,3,4\}$ be the partition supplied by global localization,
let $m_a=|\{i:a(i)=a\}|$, and put $p_a=m_a/N$.
The candidate $Z=(z_{a(i)})_i$ and its positive stress $K$ satisfy
\begin{equation}\label{app:local:stationary-data}
 ZZ^*=\frac N2 I_2,\quad K>0,\quad\tr K=2,\quad
 S_{ij}=\frac{z_i^*Kz_j}{|z_i^*Kz_j|},\quad ZS=N\lambda Z.
\end{equation}
Every stress overlap is nonzero. Put $\alpha=(1+\sqrt3)/4$.
In one common alignment with unit tetrahedral spinors $f_a$, the candidate
estimates take the form
\begin{gather}
 e=\max_a|p_a-\tfrac14|,\qquad
 \sigma^2=\sum_a(p_a-\tfrac14)^2=C/N^2,\qquad
 \|K-I\|\le\kappa,\qquad 0\le\alpha-\lambda\le\tau,
 \label{app:local:errors}\\
 \max_a\|z_a-f_a\|\le\eta,\qquad
 \sum_a p_a\|z_a-f_a\|^2\le B\sigma^2,\qquad
 \eta=\sqrt{C_{\mathrm{pt}}C}/N.                                  \label{app:local:polar-input}
\end{gather}
Here $B,C_{\mathrm{pt}}$ are the simultaneous polar bounds of
Section~\ref{app:candidate}; the separate point and aggregate estimates
must use the same alignment. The four surviving profiles are
\[
\begin{array}{c|cccc}
 C&3/4&11/4&1&3\\ \hline
 m_*=\max_a|m_a-N/4|&3/4&5/4&1/2&3/2\\
 L_*=\sum_a|m_a-N/4|&3/2&3&2&3
\end{array}
\]
with both signed profiles and all permutations included.
Throughout the certified range, $e\le1/120$ and $\eta,\kappa<1/32$.
Write $p_-=1/4-e$, $p_+=1/4+e$, and $\omega_p=1-4e>0$.

The competitor is any normalized left kernel $Y$ with
$YY^*=NI_2/2$ and $YA=0$. There is no assumption that its columns have
equal lengths or take only four values. On a closed budget region,
global localization supplies nonnegative numbers $A_R,P_R,L_U,L_P,\bar E$
such that, in its preliminary alignment,
\begin{gather}
 \|Y-F\|_F^2\le A_R/N,\qquad
 \max_i d_{\rm ph}(y_i,f_{a(i)})^2\le P_R/N,
 \label{app:local:physical-input}\\
 \sum_i d_{\rm ph}(u_i,f_{a(i)})^2\le L_U\bar E/N,\qquad
 \max_i d_{\rm ph}(u_i,f_{a(i)})^2\le L_P\bar E/N,       \label{app:local:unit-input}
\end{gather}
where $u_i=y_i/\|y_i\|$ and $d_{\rm ph}$ minimizes over a scalar phase.
The first norm is evaluated in the preliminary phase choices.
The physical lengths lie in $[1-s,1+s]$, with $s<1$.
These bounds hold throughout the region.

For fixed $K$ define
\[
 \Gamma_K(Y)=\sum_{i,j}|y_i^*Ky_j|,\qquad
 \mathcal L(Y)=\Gamma_K(Y)-N\lambda\tr(KYY^*),\qquad
 Q(H,J)=\tfrac12D^2\mathcal L(Z)[H,J].
\]
When one argument is displayed, $Q(H)$ means $Q(H,H)$.

\begin{lemma}[Closest gauge]\label{app:local:gauge}
Choose a minimizer of $\|UYD-Z\|_F$ over $U\in U(2)$ and diagonal
unitary $D$, and relabel that representative as $Y$. Set
$\Delta=Y-Z$, $d=\|\Delta\|_F$, and
\[
 C_n=\Delta\Delta^*/N,\qquad R_n=C_nZ,\qquad \Xi=\Delta+R_n .
\]
Then
\begin{gather}
 D\mathcal L(Z)=0,\qquad \Gamma_K(Z)=N^2\lambda,\qquad
 \Delta Z^*=Z\Delta^*=-\tfrac12\Delta\Delta^*,            \label{app:local:gauge-id}\\
 \Xi Z^*=0,\quad \Ima(z_i^*\Xi_i)=0,\quad
 \langle\Xi,R_n\rangle=0,\nonumber\\
 d^2=\|\Xi\|_F^2+\|R_n\|_F^2,\quad
 \|R_n\|_F/d\le d/\sqrt{2N}\quad(d>0).                 \label{app:local:normal-id}
\end{gather}
\end{lemma}
\begin{proof}
The orbit is compact. Differentiating its squared distance at a minimum
gives $\langle\Delta,AY\rangle=0$ for every skew-Hermitian $A$ and
$\langle\Delta,iY\Phi\rangle=0$ for every real diagonal $\Phi$.
The terms with $Y$ replaced by $\Delta$ vanish, so $\Delta Z^*$ is
Hermitian and $\Ima(z_i^*\Delta_i)=0$. Equal covariances give
$\Delta Z^*+Z\Delta^*=-\Delta\Delta^*$.
Multiplying the definition of $\Xi$ by $Z^*$ gives its horizontal
constraint. Since $C_n$ is Hermitian, its individual phase equations
also hold. Horizontal arrays are orthogonal to $C_nZ$, and
\[
 \|R_n\|_F^2=\frac{\tr((\Delta\Delta^*)^2)}{2N}
 \le\frac{d^4}{2N}.
\]
Finally $\nabla\Gamma_K(Z)=2KZS=2N\lambda KZ$ proves stationarity, and
$\Gamma_K(Z)=\tr(KZSZ^*)=N^2\lambda$.
\end{proof}

The simultaneous change of matrix is $A\mapsto D^*AD$:
$(UYD)(D^*AD)=UYA D=0$. Thus diagonal entries, rank and off-diagonal
moduli are preserved. We choose the representative first and then fix
its stress $K$; $\Gamma_K$ need not be invariant under common unitaries
when $K$ is fixed.
For equal covariances,
$\Gamma_K(Y)-\Gamma_K(Z)=\mathcal L(Y)-\mathcal L(Z)$.

Split $\Xi=H+U$ into cluster means and zero-sum within-cluster residuals.
This is an orthogonal splitting in the Frobenius metric.
Both $H$ and $U$ satisfy the horizontal and individual phase equations.
Every Hessian coefficient at $Z$ depends only on its two cluster types;
summing one residual factor gives zero. Consequently
\begin{equation}\label{app:local:split}
 Q(H,U)=Q(U,R_n)=0 .
\end{equation}

\subsection{The reference tensor and an endpoint perturbation bound}
\label{app:local:reference}

For a cluster-constant array, put
$\|H\|_\nu^2=\frac14\sum_a\|H_a\|^2$. More generally give each member
of cluster $a$ weight $\nu_i=1/(4m_a)$.
The normalized reference potential is
\[
 \ell_0(Y)=\sum_{i,j}\nu_i\nu_j|y_i^*y_j|
                  -\alpha\sum_i\nu_i\|y_i\|^2,\qquad
 B_0=D^2\ell_0(F),\quad q_0=B_0/2 .
\]
It has zero gradient at $F$: its phase matrix is
$\sqrt3 f_i^*f_j+(1-\sqrt3)\mathbf1_{a(i)=a(j)}$.
All infinitesimal common-unitary and individual-phase directions lie
in the full kernel of $B_0$, including their pairing with ambient vectors.

\begin{lemma}[Complete reference mean tensor]\label{app:local:reference-spectrum}
On
\[
 \mathcal H_0=\{H:\tfrac14\sum_aH_af_a^*=0,\
                         \Ima(f_a^*H_a)=0\},
\]
the full mean Hessian has eigenvalues $-2/\sqrt3$ with multiplicity three
and $-\sqrt3$ with multiplicity two. On traceless Hermitian normal
arrays $MF$, it equals $(2/\sqrt3)I$; the scalar normal is null.
The horizontal-to-normal operator norm is $\sqrt{2/3}$, and the full
ambient mean operator norm is at most $7/4$.
\end{lemma}
\begin{proof}
Let $v_a$ be tetrahedral Bloch vectors:
$\sum_av_a=0$, $\sum_av_av_a^T=4I_3/3$.
Every phase-horizontal mean has the unique representation
\[
 H_a=a_af_a+\tfrac12(u_a\cdot\sigma)f_a,\qquad
 a_a\in\mathbb R,\quad u_a\perp v_a,\qquad
 \|H_a\|^2=a_a^2+\|u_a\|^2/4.
\]
Here $\sigma$ denotes the Pauli matrices. Put
$A_1=\sum_aa_av_a$, $U_1=\sum_au_a$ and
$D_{ab}=u_a\cdot v_b+v_a\cdot u_b$.
Differentiating
$r_ar_b\sqrt{(1+x_a\cdot x_b)/2}$ at
$r_a=1$, $x_a=v_a$ gives
\begin{align}
 B_0(H,H)={}&-\frac1{2\sqrt3}\sum_aa_a^2
 +\frac{(\sum_aa_a)^2}{8\sqrt3}
 +\frac{\sqrt3}{16}A_1\cdot U_1
 +\frac{\sqrt3}{32}\|U_1\|^2
 -\frac{3\sqrt3}{256}\sum_{a\ne b}D_{ab}^2 .
 \label{app:local:reference-form}
\end{align}
For verification, along a spherical curve with initial velocity $u_a$,
the cross-pair square-root factor has derivatives
$D_{ab}/(4h)$ and
\[
 \frac{u_a\cdot u_b+(\|u_a\|^2+\|u_b\|^2)/6}{2h}
             -\frac{D_{ab}^2}{16h^3},\qquad h=1/\sqrt3.
\]
Summing ordered pairs with weight $1/16$ gives
\eqref{app:local:reference-form}; accelerations do not matter at a
stationary point.

The matrix constraint is equivalent to
$\sum a_a=0$, $U_1=-2A_1$, and $\sum u_a\times v_a=0$.
Set $b=3A_1/4$. The component
$a_a=b\cdot v_a$, $u_a=-(b-a_av_a)$ has squared $\nu$ norm
$\|b\|^2/2$ and Hessian $-\sqrt3\|b\|^2/3$.
Its orthogonal complement has $a_a=0$ and $U_1=0$.
Writing $V=(v_a)$ and $U=(u_a)$, the latter is $U=TV$ with $T$ symmetric,
and $D=2V^TTV$ gives $\|D\|_F^2=(16/3)\|U\|_F^2$.
In coordinates $v_a=(\pm1,\pm1,\pm1)/\sqrt3$ with an even number of
minus signs, $v_a^TTv_a=0$ makes $T$ diagonal and traceless.
This gives the two shape eigenvalues $-\sqrt3$.

For $H_n=(b\cdot\sigma)F$, substitution and polarization give
\[
 \|H_n\|_\nu^2=\|b\|^2,\quad
 B_0(H_n,H_n)=2\|b\|^2/\sqrt3,\quad
 B_0(H,H_n)=-\sqrt3\,b\cdot A_1/4 .
\]
The horizontal constraints imply
$\|H\|_\nu^2\ge9\|A_1\|^2/32$, with equality on the first sector.
Thus the mixed norm is $\sqrt{2/3}$. Its coupled blocks have eigenvalues
$\pm\sqrt2$, the shape eigenvalues have modulus $\sqrt3$, and all
remaining gauge/scalar modes are zero. Hence the ambient norm is
$\sqrt3<7/4$.
\end{proof}

\begin{lemma}[Squared reference projection]\label{app:local:projection}
Suppose a mean array $G$ satisfies
\[
 \left\|\tfrac14\sum_aG_af_a^*\right\|_F\le b\|G\|_\nu,
 \qquad |\Ima(f_a^*G_a)|\le\phi\|G_a\|.
\]
Its reference-horizontal and Hermitian-normal components satisfy
\[
 \|G_0\|_\nu^2\ge[1-\phi^2-3(b+\phi)^2]\|G\|_\nu^2,\qquad
 \|G_n\|_\nu\le\sqrt2 b\|G\|_\nu .
\]
Consequently, if $\beta=\phi^2+3(b+\phi)^2<1$,
\begin{equation}\label{app:local:reference-gap}
 B_0(G,G)\le-g(b,\phi)\|G\|_\nu^2,\quad
 g(b,\phi)=\frac2{\sqrt3}(1-\beta)
                         -\frac4{\sqrt3}(b+b^2).
\end{equation}
\end{lemma}
\begin{proof}
Remove the reference phase component, of norm at most $\phi\|G\|_\nu$.
The remaining constraint matrix has norm at most $(b+\phi)\|G\|_\nu$.
On phase-horizontal arrays the adjoint constraint is
$M\mapsto P_{\rm ph}(Mf_a)$, and for $M=A+iB$ with $A,B$ Hermitian,
$\tr B=0$,
\[
 \|P_{\rm ph}(MF)\|_\nu^2=\tfrac12\|A\|_F^2+\tfrac13\|B\|_F^2 .
\]
Thus its least nonzero singular value is $1/\sqrt3$.
The two successive corrections are orthogonal, proving the squared
horizontal estimate. If $C=\frac14\sum G_af_a^*$, the normal projection
is $2\operatorname{Herm}(C)F$ and has norm squared
$2\|\operatorname{Herm}(C)\|_F^2$. Reference gauge directions are
orthogonal to this normal and annihilated by $B_0$.
Lemma~\ref{app:local:reference-spectrum}, $\|G_0\|_\nu\le\|G\|_\nu$,
and the normal estimate yield \eqref{app:local:reference-gap}.
\end{proof}

The following endpoint version of the modulus derivative estimate is
useful because a certified overlap floor concerns actual overlaps,
not the intervening line segment.
\begin{lemma}[Endpoint Hessian bound]\label{app:local:endpoint-hessian}
For nonzero $u,v\in\mathbb C\simeq\mathbb R^2$, and
$0<h\le\min(|u|,|v|)$,
\[
 \|D|\cdot|(u)-D|\cdot|(v)\|\le |u-v|/h,\qquad
 \|D^2|\cdot|(u)-D^2|\cdot|(v)\|
                           \le\frac76\,|u-v|/h^2 .
\]
\end{lemma}
\begin{proof}
The case $u=v$ is immediate. Write $r=|u|\ge s=|v|$,
$t=r/s$, $d=t-1$, $c=\cos\arg(u\bar v)$, and
$D=|u-v|/s=\sqrt{d^2+2t(1-c)}$.
The gradient difference has squared norm $(D^2-d^2)/t\le D^2$.
The norm of the Hessian difference, multiplied by $s$, is exactly
\[
 \frac{d+\sqrt{d^2+4t(1-c^2)}}{2t}.
\]
Put $L=7/6$ and $y=d/D$. Since
$1-c^2\le(D^2-d^2)/t$,
\[
 tL^2D^2-dLD-(1-c^2)
 \ge\frac{D^2}{t}\left[(y-tL/2)^2+\frac34t^2L^2-1\right]>0.
\]
Also $2tLD\ge d$. Squaring in the correct direction proves the displayed
Hessian norm is at most $LD$. Rescaling and $s\ge h$ prove the claim.
No lower bound on intermediate segment moduli was assumed.
\end{proof}

For later use define, for a certified cross stress floor $0<h_c\le15/26$,
\begin{align}
 d_p&=2\eta+\eta^2+\kappa(1+\eta)^2,\quad
 u_p=\eta+\kappa(1+\eta), \notag\\
 T_p&=\frac{7d_p}{6h_c^2}
       +\frac{u_p(2+u_p)}{h_c}+\frac{d_p}{2h_c}+\frac\kappa2 .
 \label{app:local:pair-error}
\end{align}
Then the absolute \emph{full} cross-pair Hessian error is at most
$T_p(\|h_a\|+\|h_b\|)^2$. Indeed the overlap, first-variation and
second-variation errors are bounded respectively by $d_p$, $u_pv$,
and $\kappa v^2/2$, where $v=\|h_a\|+\|h_b\|$; the corresponding reference
variations have bounds $v$ and $v^2/2$.
Apply Lemma~\ref{app:local:endpoint-hessian} to the reference
first-variation arguments, then change those arguments, the gradient,
and the second variation. This gives precisely the four terms in
\eqref{app:local:pair-error}. Same-type \emph{mean} pairs have exact
Hessian $2h_a^*Kh_a$, and hence error at most $2\kappa\|h_a\|^2$.

\subsection{Mean curvature in physical, mass, and stress coordinates}
\label{app:local:mean}

An ambient physical mean perturbation bound needed for the normal terms is
\begin{align}
 W_\times&=\max_{a\ne b}|16p_ap_b-1|,\qquad W_d=8e+16e^2,\notag\\
 E_{\mathrm{Hess}}&=(1+4e)^2(3T_p+\kappa/2)
       +(27/4)W_\times+W_d/2 \notag\\
 &\quad+2\{(1+4e)[(7/10+\tau)\kappa+\tau]+(14/5)e\}.
 \label{app:local:ambient-P}
\end{align}
It satisfies
\begin{equation}\label{app:local:ambient-perturbation}
 |N^{-2}Q(H,J)-q_0(H,J)|
                  \le(E_{\mathrm{Hess}}/2)\|H\|_\nu\|J\|_\nu
\end{equation}
on all cluster-constant arrays, without a horizontal hypothesis.
To verify this, use
$\sum_{a\ne b}(l_a+l_b)^2/16\le3\|H\|_\nu^2$.
The reference cross Hessian is bounded by $(9/4)(l_a+l_b)^2$;
the same-type one is exactly $2l_a^2$. This gives the cross and diagonal
weight errors separately. The multiplier error follows from
$\|\lambda K-\alpha I\|\le(7/10+\tau)\kappa+\tau$ and
$|4p_a-1|\le4e$. Polarization of the resulting absolute self-adjoint
operator bound gives \eqref{app:local:ambient-perturbation}.
For the four profiles, respectively, valid cross-weight envelopes are
\[
 W_\times\le 2/N+3/N^2,\quad 6/N+15/N^2,\quad
                 4/N+4/N^2,\quad 4/N+12/N^2 .
\]
They follow directly from
$16p_ap_b-1=4(\delta_a+\delta_b)/N+16\delta_a\delta_b/N^2$.

\begin{lemma}[Mass coordinates]\label{app:local:mass-mean}
Let $\epsilon_+\ge\max_a(p_a-1/4)$ cover both signs of the count profile,
and define
\[
 b_g=\sqrt{B\sigma^2/\omega_p},\quad
 g_g=g(b_g,\eta)-(1+2/\omega_p)T_p-\kappa/2
                   -2(\alpha\kappa+\tau)/\omega_p-8\alpha\epsilon_+/\omega_p .
\]
If the reference projection fraction is positive and $g_g>0$, then
\begin{equation}\label{app:local:mass-gap}
 Q(H)\le-Nc_g\|H\|_F^2,\qquad c_g=\omega_p g_g/2
\end{equation}
for every actual horizontal mean.
\end{lemma}
\begin{proof}
Set $G_a=4p_aH_a$. Then $\frac14\sum G_az_a^*=0$ and
$\Ima(z_a^*G_a)=0$. The reference constraint error is at most
$b_g\|G\|_\nu$, and its phase error is at most $\eta$.
Lemma~\ref{app:local:projection} gives $g(b_g,\eta)$.
The exact norm relation is
$\|H\|_F^2/N=\frac14\sum\|G_a\|^2/(4p_a)\le\|G\|_\nu^2/\omega_p$.

For completeness, the mass cancellation can be checked at the
half-Hessian level. Put
$\theta^0_{ab}=f_a^*f_b/|f_a^*f_b|$ and
$A_{ab}(g)=[\Ima(\overline{\theta^0_{ab}}g^*f_b)]^2/|f_a^*f_b|$.
If $q_p^F$ denotes the reference half-Hessian with masses $p$, then
\[
 q_p^F(H)-q_{1/4}^F(G)
 =\sum_a\frac{\sum_b(p_b-p_a)A_{ab}(G_a)
                         +4\alpha(p_a-\tfrac14)\|G_a\|^2}{4(4p_a)} .
\]
Mixed first-variation and second-variation terms cancel because
$p_ap_b/[(4p_a)(4p_b)]=1/16$.
In coordinates $f_a=(1,0)$, $g=(x+iy,u+iv)$,
$\sum_bA_{ab}(g)=(1+\sqrt3)y^2+\sqrt3(u^2+v^2)
\le4\alpha\|g\|^2$.
The displayed difference is therefore at most
$4\alpha\epsilon_+\|G\|_\nu^2/\omega_p$.

Moreover
\[
 \sum_{a\ne b}p_ap_b(\|H_a\|+\|H_b\|)^2
                       \le(1+2/\omega_p)\|G\|_\nu^2 .
\]
Expanding the two squared-column terms gives at most
$(4/\omega_p-1)\sum\|G_a\|^2/8$, and the product terms at most
$3\sum\|G_a\|^2/8$.
The same-type full-Hessian parameter error is $\kappa\|G\|_\nu^2/2$,
and the covariance error is $2(\alpha\kappa+\tau)\|G\|_\nu^2/\omega_p$.
These estimates give $g_g$. Its strict positivity is checked before
converting the negative norm estimate.
\end{proof}

\begin{lemma}[Stress and mass coordinates]\label{app:local:stress-mean}
Put
\begin{gather*}
 d_K=\frac{\kappa}{1+\sqrt{1-\kappa}},\qquad
 \eta_v=\sqrt{1+\kappa}\eta+d_K,\qquad
 b_v=\frac{\sqrt{B\sigma^2}+d_K}{\sqrt{\omega_p}},\\
 \phi_v=\eta_v+\frac{\kappa(1+\eta)}{\sqrt{1-\kappa}},\qquad
 d_v=2\eta_v+\eta_v^2,\\
 T_v=\frac{7d_v}{6h_c^2}
               +\frac{\eta_v(2+\eta_v)}{h_c}+\frac{d_v}{2h_c},\\
 g_v=g(b_v,\phi_v)-(1+2/\omega_p)T_v-2\tau/\omega_p-8\alpha\epsilon_+/\omega_p .
\end{gather*}
Assume $\eta_v<1/32$, $\phi_v^2+3(b_v+\phi_v)^2<1$ and $g_v>0$.
Then
\begin{equation}\label{app:local:stress-gap}
 Q(H)\le-Nc_v\|H\|_F^2,\qquad c_v=\omega_p(1-\kappa)g_v/2 .
\end{equation}
\end{lemma}
\begin{proof}
Set $T=K^{1/2}$, $v_a=Tz_a$, and $G_a=4p_aTH_a$.
The entire Lagrangian in $TY$ has pair stress $I$ and multiplier
$\lambda I$. The transformed covariance is $K/2$, not $I/2$;
no tightness of $v$ is used.
Exactly $\frac14\sum G_av_a^*=0$, but its phase equation has the defect
\[
 |\Ima(v_a^*G_a)|
 =4p_a|\Ima(z_a^*(K-I)H_a)|
 \le\frac{\kappa(1+\eta)}{\sqrt{1-\kappa}}\|G_a\|.
\]
Spectral calculus gives $\|T-I\|\le d_K$, proving the point bound
$\|v_a-f_a\|\le\eta_v$. For the aggregate use
$v-f=(z-f)+(T-I)z$ and the \emph{original} identity
$\sum p_a\|z_a\|^2=1$:
\[
 \left(\sum p_a\|v_a-f_a\|^2\right)^{1/2}
                         \le\sqrt{B\sigma^2}+d_K .
\]
Thus Lemma~\ref{app:local:projection} applies with $b_v,\phi_v$;
in particular the transformed phase defect has not been erased.

The cross-pair perturbation is now
\eqref{app:local:pair-error} with point error $\eta_v$ and stress
error zero, namely $T_v$. The overlap floor is unchanged, since
$v_a^*v_b=z_a^*Kz_b$.
Same-type parameter error is exactly zero, because both kernels have
Hessian $2\|g_a\|^2$. Their mass error is still included in the preceding
mass-coordinate identity. The covariance error is $2\tau/\omega_p$.
This proves the Hessian bound with coefficient $g_v$.
Finally
\[
 \|G\|_\nu^2=\sum_a4p_a^2H_a^*KH_a
             \ge \omega_p(1-\kappa)\|H\|_F^2/N .
\]
Its substitution into a negative bound is valid after $g_v>0$.
The transformation is invertible and introduces no equality mode.
\end{proof}

The accepted mean coefficient may be any certified lower bound for
$\max(c_g,c_v)$, using only those alternatives whose guards hold.
The stress transform is confined to this mean estimate. In particular
neither $b_g$ nor $b_v$ replaces the physical normal projection below.

\subsection{The intrinsic residual bound}
\label{app:local:residual}

We use the two-dimensional spectral reduction, since a trace-only
bound would lose the anisotropy needed in the nonuniform profiles.
Use the candidate construction of Section~\ref{app:candidate}, writing
its auxiliary rank-two Gram matrix as $\mathcal B$ to distinguish it
from the aggregate coefficient $B$ in \eqref{app:local:polar-input}.
Define
\begin{gather}
 \ell_a=\frac{1-\lambda-p_a}{(1-\lambda)(\lambda-2p_a)},\quad
 L=\diag(\ell_a),\quad \mathcal T=L-L\mathcal B L,\quad
 \mathcal T_0=\tr\mathcal T, \notag\\
 d_a=\ell_a[1-\ell_a(\lambda-p_a)],\quad
 r_{Ka}^2=\frac{d_a}{p_a\mathcal T_0},\qquad
 s_{ab}=\frac{\ell_a\ell_b}{\mathcal T_0}\quad(a\ne b),\quad
 s_{aa}=r_{Ka}^2.                                    \label{app:local:intrinsic-scalars}
\end{gather}
All displayed denominators and $d_a$ are strictly positive.
If the columns of $U$ form an orthonormal basis of
$\operatorname{range}\mathcal T$, then
$K=2U^*\mathcal T U/\mathcal T_0$ and
$z_a=U^*e_a/\sqrt{2p_a}$. The exact identity
$\mathcal T L^{-1}\mathcal T=\mathcal T$ implies, for $v_a=K^{1/2}z_a$,
\begin{equation}\label{app:local:weighted-resolution}
 \sum_a(p_a/\ell_a)v_av_a^*=I_2/\mathcal T_0 .
\end{equation}
Indeed $U^*L^{-1}U=(U^*\mathcal T U)^{-1}$.

\begin{lemma}[Pointwise intrinsic form]\label{app:local:intrinsic}
For a physical vector $h$ at type $a$ define
\[
 q_a(h)=\sum_b\frac{p_b}{s_{ba}}
       [\Ima(\overline{\theta_{ba}}z_b^*Kh)]^2-\lambda h^*Kh,
 \qquad \theta_{ba}=z_b^*Kz_a/s_{ba}.
\]
For the full within-cluster zero-sum residual $U$,
$Q(U)/N=\sum_jq_{a(j)}(U_j)$.
For $j\ne a$ put
\[
 u_j=p_j^2(r_{Kj}^2-s_{aj}^2/r_{Ka}^2),\qquad
 D_a=\frac{2\sum_{j<k}u_ju_k-\sum_ju_j^2}{4},
\]
where the sums run over the three indices different from $a$, and set
\begin{gather}
 \mathcal D_a=D_a\sum_{\substack{j<k\\j,k\ne a}}
                       \frac1{p_jp_ks_{aj}s_{ak}},\qquad
 c_{Ka}=\lambda-\frac{\ell_a^{-1}
                   +\sqrt{\ell_a^{-2}-4\mathcal D_a}}2 .
 \label{app:local:intrinsic-eigen}
\end{gather}
If $c_K\le\min_a\{\lambda,c_{Ka}\}$ is strictly positive, then
\begin{equation}\label{app:local:physical-residual}
 q_a(h)\le-(1-\kappa)c_K\|h\|^2
 \quad\text{whenever }\Ima(z_a^*h)=0 .
\end{equation}
\end{lemma}
\begin{proof}
In the ordered Hessian, mixed-argument and mixed first-variation terms
of a zero-sum residual factor through its group sums and vanish.
The two pure-first terms give $Q(U)/N=\sum_jq_{a(j)}(U_j)$,
including same-cluster and diagonal pairs.

Fix $e_0=v_a/r_{Ka}$ and a perpendicular unit vector $e_1$.
Rotate $v_b$ into its pair phase, writing
$(s_{ab}/r_{Ka})e_0+c_be_1$, with $c_a=0$.
Stress stationarity gives
$\sum_bp_bc_b=0$ and $\sum_bp_bs_{ab}=\lambda r_{Ka}^2$.
If $K^{1/2}h=(x+iy)e_0+u e_1$, direct expansion therefore yields
\[
 q_a(h)=-\lambda x^2+
         \sum_{b\ne a}\frac{p_b}{s_{ab}}
                [\Ima(\bar u c_b)]^2-\lambda|u|^2 .
\]
All terms containing $y$ cancel exactly; the whole $iz_a$ direction
is in the kernel, including its bilinear pairings.

Let $M_a$ be the real two-dimensional positive quadratic form in the
middle sum. Each real rank-one form $[\Ima(\bar u c_b)]^2$
has trace $|c_b|^2$. Since $c_a=0$ and $s_{ab}=\ell_a\ell_b/\mathcal T_0$ for $b\ne a$,
\begin{align*}
 \tr M_a&=\sum_{b\ne a}\frac{p_b|c_b|^2}{s_{ab}}
 =\frac{\mathcal T_0}{\ell_a}\sum_b\frac{p_b}{\ell_b}|c_b|^2,\\
 \sum_b\frac{p_b}{\ell_b}|c_b|^2
 &=e_1^*\left(\sum_b\frac{p_b}{\ell_b}v_bv_b^*\right)e_1
 =\frac1{\mathcal T_0}.
\end{align*}
The last equality is \eqref{app:local:weighted-resolution}; pair-phase
rotations do not change the projected moduli. Thus $\tr M_a=1/\ell_a$.
The vectors $t_j=p_jc_j$ form a closed triangle. Their pairwise
determinants have common square $D_a$, since
$t_j\cdot t_k=(u_l-u_j-u_k)/2$.
The two-dimensional Cauchy--Binet formula gives $\det M_a=\mathcal D_a$.
This proves its \emph{complete} largest-eigenvalue formula
\eqref{app:local:intrinsic-eigen}.

For a physical phase-horizontal $h$, choose real $t$ so that
$h'=h-it z_a$ satisfies $\Ima(z_a^*Kh')=0$.
Then $q_a(h')=q_a(h)$, while
$\|h'\|^2=\|h\|^2+t^2\|z_a\|^2\ge\|h\|^2$.
The spectral bound and $K\ge(1-\kappa)I$ prove
\eqref{app:local:physical-residual}.
This is a pointwise statement, prior to imposing any zero group sum.
The phase adjustment is real-linear within each group, so it would also
preserve such sums if applied to an array.
\end{proof}

In the accepted comparisons we use a directed lower bound
\begin{equation}\label{app:local:cr}
 c_r=(1-\kappa)c_K>0
\end{equation}
from this intrinsic lemma. All residual lower bounds used in the
certificates follow from this lemma.
The scalar enclosures cover the \emph{whole} reciprocal-$N$ interval,
not merely the local endpoint. They are produced by
\eqref{app:local:intrinsic-scalars}--\eqref{app:local:intrinsic-eigen}
with the root enclosure and contractor described next.

Write $x=1/N$, $p_a=1/4+\delta_ax$, $h_0=\alpha-1/2$,
$\lambda=\alpha-\zeta$, and $h=h_0-\zeta$.
Polynomial division of the defining candidate equation gives exactly
\begin{equation}\label{app:local:root-contractor}
 \zeta=x^2
 \frac{3+1/h-1/(4h^2)}{2+1/(4h_0h)}
                 \sum_a\frac{\delta_a^2}{h-2\delta_ax}.
\end{equation}
For verification,
$\sum_a(h-2\delta_ax)^{-1}
=4/h+(4x^2/h^2)\sum_a\delta_a^2/(h-2\delta_ax)$;
the uniform equation is $2h_0+1-1/(4h_0)=0$.
On any closed $x$ interval start with
$0\le\zeta\le(2/3)Cx_{\rm hi}^2$ and intersect this interval three times
with the outward interval image of \eqref{app:local:root-contractor}.
The true root remains inside each intersection; convergence is not
assumed. Positivity of every denominator is checked before division.
For the actual candidate the discriminant is nonnegative. Its enclosing
interval may therefore be intersected with $[0,\infty)$; a negative lower
endpoint is replaced by zero, while a negative upper endpoint fails the check.
Taking the minimum over all intervals, signs and four vertex types gives
$c_K$. All intervals meeting a requested endpoint are retained whole,
including any interval crossing it.

\subsection{The normal component and its finite budget}
\label{app:local:normal}

The candidate diagonal identity is
$p_a\|z_a\|^2=1/4+(p_a-1/4)/b$ with
$b=3(2\lambda-1)>109/100$. It implies
\[
 \tfrac14\sum_a\|f_a-4p_az_a\|^2
 =\sum_ap_a\|z_a-f_a\|^2+4\sigma^2/b .
\]
Set
\[
 b_{\rm ph}=\sqrt{(B+400/109)\sigma^2}.
\]
For a physical horizontal mean $H$ this bounds its reference constraint
matrix by $b_{\rm ph}\|H\|_\nu$; in particular its reference normal
projection has norm at most $\sqrt2b_{\rm ph}\|H\|_\nu$.
The following lemma uses these bounds in the original frame variables.

\begin{lemma}[Positive normal coupling]\label{app:local:PSD-normal}
Define
\[
 t=\min\{\sqrt2\eta,\sqrt{2B\sigma^2/\omega_p}\},\qquad
 a=\min\{\omega_p^{-1/2},1+t\},
\]
and
\begin{align}
 J_*&=\frac{(1/2+b_{\rm ph})/\sqrt3+(7/8)t+(E_{\mathrm{Hess}}/2)a}{\sqrt{\omega_p}},
 \notag\\
 L_*'&=\frac1{2\sqrt3}+\frac74t+\frac78t^2+(E_{\mathrm{Hess}}/2)a^2,
 \notag\\
 H_{\rm amb}&=(7/4+E_{\mathrm{Hess}})/(2\omega_p),\qquad
 J=\min(J_*,H_{\rm amb}),\quad L=\min(L_*',H_{\rm amb}).
 \label{app:local:JL}
\end{align}
Then
$|Q(H,R_n)|\le NJ\|H\|_F\|R_n\|_F$ and
$Q(R_n)\le NL\|R_n\|_F^2$.
\end{lemma}
\begin{proof}
Write $C_n=cI+D$ with $\tr D=0$ and set
$r=\|R_n\|_F/\sqrt N$. Tightness gives
$r^2=\tr(C_n^2)/2=\|C_nF\|_\nu^2$.
Because $C_n\succeq0$ in dimension two,
$\|DF\|_\nu^2=\tr(D^2)/2\le r^2/2$.
This half-energy estimate would fail for an indefinite normal.
Also $\|C_n\|\le\sqrt2r$, so for $E=C_n(Z-F)$,
$\|E\|_\nu\le tr$, and $\|R_n\|_\nu\le ar$.
Both the point and aggregate alternatives in $t$ refer to this same $E$.

By the reference tensor and the physical normal projection,
\[
 |q_0(H,C_nF)|
 \le\frac{1/2+b_{\rm ph}}{\sqrt3}\|H\|_\nu r .
\]
The error $q_0(H,E)$ is at most $(7/8)t\|H\|_\nu r$.
Add \eqref{app:local:ambient-perturbation} and use
$\|H\|_\nu\le\|H\|_F/\sqrt{N\omega_p}$ to obtain $J_*$.
For the pure normal, the scalar part is null and the traceless part
contributes at most $r^2/(2\sqrt3)$ to $q_0(C_nF)$.
The two terms involving $E$ are bounded by
$(7t/4+7t^2/8)r^2$, and the perturbation contributes
$E_{\mathrm{Hess}} a^2r^2/2$. This gives $L_*'$.
The ambient operator bound independently gives $H_{\rm amb}$ for each.
\end{proof}

Let $c_m>0$ be the accepted mean lower bound and $c_r>0$ the intrinsic
residual lower bound. Put
$h=\|H\|_F/d$, $r=\|U\|_F/d$, $z=\|R_n\|_F/d$.
Although $\Delta=\Xi-R_n$, an absolute mixed upper bound gives
\begin{equation}\label{app:local:normal-optimization}
 \frac{Q(\Delta)}{Nd^2}
 \le-c_mh^2-c_rr^2+2Jzh+Lz^2,\qquad
 h^2+r^2+z^2=1,\quad 0\le z\le\rho_n .
\end{equation}
If $d\le d_c/\sqrt N$, take $\rho_n^2=d_c^2/(2N^2)\le1/2$.
The exact supremum of the right-hand side is
\begin{equation}\label{app:local:normal-max}
 \begin{cases}
 -c_r+(c_r+L+J^2/(c_m-c_r))\rho_n^2,
 &c_m>c_r,\ J^2\rho_n^2\le(c_m-c_r)^2(1-\rho_n^2),\\
 -c_m+(c_m+L)\rho_n^2+2J\sqrt{\rho_n^2(1-\rho_n^2)},
 &\text{otherwise}.
 \end{cases}
\end{equation}
To prove this, fix $z$ and maximize
$-(c_m-c_r)h^2+2Jzh$ on $[0,\sqrt{1-z^2}]$.
The interior value and the boundary value are nondecreasing in $z$
for $z\le1/\sqrt2$, and agree at the transition.
The boundary case includes $c_m\le c_r$, so no division by a nonpositive
gap occurs. Denote the negative of \eqref{app:local:normal-max} by
$g_{\rm normal}$. All branch choices are made with the conservative
rational coefficients actually supplied to the certificate.

\subsection{Angular transport and the shared radial--tangent budget}
\label{app:local:angular}

The columnwise tangent used in the Taylor remainder is different from
the horizontal array $\Xi$ used in the Hessian. Define
\[
 D_r=\frac{(4/11)e}{1-4e},\qquad r_-=\sqrt{1-D_r},\quad
 r_+=\sqrt{1+D_r},\qquad d_c=\sqrt{A_R}+\sqrt{BC}.
\]
The candidate radial identity in Section~\ref{app:candidate} gives
$r_-\le\|z_a\|\le r_+$. Closest alignment and the triangle inequality
give $d\le d_c/\sqrt N$.
In the preliminary common alignment, the exact identity
\[
 \|rv-f\|^2=(r-1)^2+r\,d_{\rm ph}(v,f)^2
\]
for unit $v,f$ in closest phases gives candidate angular bounds
$BC/(Nr_-)$ in sum and $C_{\mathrm{pt}}C/(N^2r_-)$ pointwise. Put
\[
 B_{\rm ang}=\sqrt{L_U\bar E}+\sqrt{BC/r_-},\qquad
 \chi=\sqrt{\frac{(1+s)r_+}{(1-s)r_-}},\qquad
 D_U=\frac23-\frac{L_*}{N}>0.
\]
For fixed common unitary, minimizing individual phases gives
\[
 \|y_i-z_i\|^2=(\|y_i\|-\|z_i\|)^2+
 \|y_i\|\|z_i\|d_{\rm ph}(u_i,z_i/\|z_i\|)^2.
\]
Therefore the actual closest unitary minimizes the angular sum with
the positive weights $\|y_i\|\|z_i\|$. Comparison to the preliminary
alignment yields
\begin{align}
 \sum_i\|y_i\|\|z_i\|d_{\rm ph}(u_i,z_i/\|z_i\|)^2
 &\le (1+s)r_+B_{\rm ang}^2/N,\label{app:local:weighted-angle}\\
 \sum_i d_{\rm ph}(u_i,z_i/\|z_i\|)^2
 &\le\chi^2B_{\rm ang}^2/N. \nonumber
\end{align}
No unweighted angular optimality is used.

For completeness, the rotation estimate needed to transfer the pointwise
bound follows from the tetrahedral second moment. Modulo scalar phase
write $U=\cos\theta I+i\sin\theta(n\cdot\sigma)$ and put $t=|\sin\theta|$.
If $v_a$ are the reference Bloch vectors, then
$|f_a^*Uf_a|^2=1-t^2(1-(n\cdot v_a)^2)$ and
$d_{\rm ph}(Uf_a,f_a)^2\ge t^2(1-(n\cdot v_a)^2)$.
Since $\sum v_av_a^T=4I/3$,
\[
 \sum_i d_{\rm ph}(Uf_{a(i)},f_{a(i)})^2\ge ND_Ut^2.
\]
The phase-distance triangle inequality and
\eqref{app:local:weighted-angle} bound its left side by
$(1+\chi)^2B_{\rm ang}^2/N$. If
\begin{equation}\label{app:local:rotation-guard}
 \frac{(1+\chi)^2B_{\rm ang}^2}{N^2D_U}
 <1-(255/256)^2,
\end{equation}
every $|f_a^*Uf_a|>255/256$, and
\[
 d_{\rm ph}(Uf_a,f_a)^2=
 \frac{2(1-|f_a^*Uf_a|^2)}{1+|f_a^*Uf_a|}
 \le\frac{512}{511}t^2.
\]
Thus one may use
\[
 \mathrm{Rot}_{\rm ang}=
 \sqrt{\frac{512(1+\chi)^2B_{\rm ang}^2}{511D_U}}.
\]
An independent physical-distance triangle gives
$\sum_i d_{\rm ph}(Uf_a,f_a)^2\le4d_c^2/N$ and hence
$\mathrm{Rot}_{\rm phys}=\sqrt{(512/511)4d_c^2/D_U}$,
provided $4d_c^2/(N^2D_U)<1-(255/256)^2$.
Use the smaller of these two bounds only when each selected bound has
its own stated guard; denote the result by $\mathrm{Rot}$.

Write, in the actual individual phases,
\begin{equation}\label{app:local:column-split}
 y_i=(1+t_i)z_i+\xi_i=r_i z_i+\xi_i,\qquad
 z_i^*\xi_i=0,\quad t_i\in\mathbb R.
\end{equation}
Let $X^2=\sum_i\|\xi_i\|^2$ and let $T$ bound $\max_i\|\xi_i\|$.
The following simultaneous choices are valid:
\begin{align}
 q_t&=\frac1N\min\{d_c,(1+s)\sqrt{r_+/r_-}B_{\rm ang}\},
 & X/\sqrt N&\le q_t,\label{app:local:tangent-bounds}\\
 \epsilon&=\sqrt{P_R/N}+(\mathrm{Rot}_{\rm phys}+\sqrt{C_{\mathrm{pt}}C})/N,
 & q_d&=d_c/N,\nonumber\\
 T&=\min\left\{\epsilon,(1+s)\left[
 \sqrt{L_P\bar E/N}+\frac{\mathrm{Rot}}N+
 \frac{\sqrt{C_{\mathrm{pt}}C}}{N\sqrt{r_-}}\right]\right\}.&&\nonumber
\end{align}
Indeed $1-|u^*v|^2\le d_{\rm ph}(u,v)^2$ for unit vectors.
Multiplying \eqref{app:local:weighted-angle} by $(1+s)/r_-$ gives
$X^2\le(1+s)^2(r_+/r_-)B_{\rm ang}^2/N$.
The pointwise estimate follows from the same triangle inequality with
the rotation bound, while the physical triangle gives $\epsilon$.
An angular displacement has therefore not been charged a first-order
radial point error.

The certificate checks $T<1-s$. Since $z_i^*y_i\ge0$ in closest phases,
\begin{equation}\label{app:local:radial-interval}
 0<\rho=\frac{\sqrt{(1-s)^2-T^2}}{r_+}\le r_i
 \le R=\frac{1+s}{r_-}.
\end{equation}
Orthogonality in \eqref{app:local:column-split} gives the exact identity
\begin{equation}\label{app:local:shared-budget}
 d^2=\sum_i t_i^2\|z_i\|^2+X^2,\qquad
 T_r^2:=\sum_i t_i^2\le(d^2-X^2)/r_-^2.
\end{equation}
This budget is distinct from the global radial, moment and angular
deficits. Equations \eqref{app:local:tangent-bounds} are the explicit
transfer between them. For $d>0$ put $\vartheta=X/d\in[0,1]$.

\subsection{The pair remainder and physical fiber bounds}
\label{app:local:pair-remainder}

For an ordered pair put $s_{ij}=|z_i^*Kz_j|$, and rotate its candidate
overlap to the positive real axis. In this phase set
\[
 \alpha_{ij}=\overline S_{ij}z_i^*K\xi_j,\quad
 \beta_{ij}=\overline S_{ij}\xi_i^*Kz_j,\quad
 \gamma_{ij}=\overline S_{ij}\xi_i^*K\xi_j,
 \qquad (a_{ij},b_{ij},c_{ij})=\Ima(\alpha_{ij},\beta_{ij},\gamma_{ij}).
\]
For this paragraph abbreviate $s=s_{ij}$, $p=r_ir_j$ and
$L=r_i\alpha_{ij}+r_j\beta_{ij}$. The complete Taylor polynomial
through total degree two in $(t,\xi)$ is
\[
 P_2=ps+\Rea(L+\gamma_{ij})+(a_{ij}+b_{ij})^2/(2s).
\]
With $R_s(h)=|s+h|-s-\Rea h-(\Ima h)^2/(2s)$, the exact remainder is
\begin{align}
 |ps+L+\gamma_{ij}|-P_2={}&
 \frac{r_i-r_j}{2s}\left(\frac{a_{ij}^2}{r_j}-\frac{b_{ij}^2}{r_i}\right)
 +\frac{(r_i a_{ij}+r_j b_{ij})c_{ij}}{ps}
 +\frac{c_{ij}^2}{2ps}\nonumber\\
 &+pR_s\left(\frac{\alpha_{ij}}{r_j}+
              \frac{\beta_{ij}}{r_i}+\frac{\gamma_{ij}}p\right).
 \label{app:local:exact-pair}
\end{align}
To verify it, expand the square of
$(r_i a_{ij}+r_j b_{ij}+c_{ij})$ divided by $2ps$;
the $a_{ij}b_{ij}$ term cancels exactly. In particular the pure radial
dependence is exactly quadratic. All within-cluster and diagonal pairs
remain in this identity.

For types $\mathrm c$ (different clusters) and $\mathrm s$ (same cluster),
choose positive floors $h_t$ and upper bounds $m_t,g_t$ such that
\begin{equation}\label{app:local:type-bounds}
 s_{ij}\ge h_t,\quad |\alpha_{ij}|\le m_t\|\xi_j\|,\quad
 |\beta_{ij}|\le m_t\|\xi_i\|,\quad
 |\gamma_{ij}|\le g_t\|\xi_i\|\|\xi_j\|.
\end{equation}
We use $h_{\mathrm c}=h_c$ for the certified cross floor.
Besides the certified candidate overlap floor, valid generic choices are
\begin{gather}
 h_E=15/26-2\eta-\eta^2>0,\qquad
 h_{\mathrm s}=(1-\kappa)r_-^2,\nonumber\\
 m_{\mathrm c}=\sqrt{r_+^2-h_E^2/r_+^2}+\kappa r_+,\qquad
 m_{\mathrm s}=\kappa r_+,\label{app:local:generic-fiber}\\
 g_{\mathrm c}=\kappa+\min\{1,(1/\sqrt3+2\eta+\eta^2)/r_-^2\},\qquad
 g_{\mathrm s}=1+\kappa.\nonumber
\end{gather}
The displayed choice of $h_{\mathrm s}$ may be replaced by another
independently proved floor; in all cases require
$0<h_{\mathrm s}\le\min_a s_{aa}$.
The projected-linear bound follows from the Euclidean cross floor and
the component of $K-I$. For the quadratic bound, define the antiunitary
$J(v_1,v_2)=(-\overline v_2,\overline v_1)$. Each
$\xi_i=c_iJ(z_i/\|z_i\|)$, so
\[
 |\xi_i^*\xi_j|=
 \frac{|z_i^*z_j|}{\|z_i\|\|z_j\|}\|\xi_i\|\|\xi_j\|.
\]
Adding the $K-I$ contribution proves $g_{\mathrm c},g_{\mathrm s}$.
In particular this is a physical complex-line calculation, valid for
every within-cluster residual.

\subsection{Row estimates coupled to intrinsic residual curvature}
\label{app:local:rows}

Write $e_i=\|\xi_i\|$. Retain the separate pointwise coefficient
\[
 M=\max\{m_{\mathrm c}^2/h_{\mathrm c},m_{\mathrm s}^2/h_{\mathrm s}\}.
\]
The following row bounds keep the complete imaginary parts in
\eqref{app:local:exact-pair}:
\[
 \sum_i a_{ij}^2/s_{ij}\le NH e_j^2,\qquad
 \sum_{a(i)\ne a(j)}a_{ij}^2/s_{ij}\le NH_{\mathrm c}e_j^2.
\]
To obtain an initial bound, put
\begin{gather}
 d_g=2\eta+\eta^2+\kappa r_+^2,\qquad
 \delta=\eta+\kappa r_++\eta/\sqrt{r_-}
       +\sqrt{2/3}\,d_g/\sqrt{h_{\mathrm c}/\sqrt3},\nonumber\\
 H_{\rm iso}=\frac{(\sqrt{p_+}+\sqrt{1-p_-}\,\delta)^2}{h_{\mathrm c}},
 \label{app:local:row-initial}\\
 H=\min\{M,H_{\rm iso}+p_+m_{\mathrm s}^2/h_{\mathrm s}\},\qquad
 H_{\mathrm c}=\min\{H_{\rm iso},(1-p_-)m_{\mathrm c}^2/h_{\mathrm c}\}.
 \nonumber
\end{gather}
The subscript on $H_{\rm iso}$ denotes the reference isotropy estimate.
For a reference center $f_b$, the three projected cross forms have
amplitude $\sqrt{2/3}\|\xi\|$ and phases separated by $2\pi/3$.
Their squared imaginary parts sum to $\|\xi\|^2$, so their weighted sum
is at most $p_+\|\xi\|^2$. Transport $\xi=cJ(z_b/\|z_b\|)$ to
$cJf_b$; its change is at most $\eta\|\xi\|/\sqrt{r_-}$.
For complex nonzero $v,v_0$,
\[
 \left|\frac v{|v|}-\frac{v_0}{|v_0|}\right|^2
 =\frac{|v-v_0|^2-(|v|-|v_0|)^2}{|v||v_0|}.
\]
The overlap phase changes by at most $d_g/\sqrt{h_{\mathrm c}/\sqrt3}$.
These bounds give the displayed $\delta$ for the full phase-aligned
linear form. The weighted Euclidean triangle inequality proves
$H_{\rm iso}$; same-cluster rows and pointwise $M$ give the other terms.

There is a second, sharper row bound directly from intrinsic curvature.
Let $c_{\rm intr}$ be the physical pointwise coefficient established in
the residual argument above, before imposing any zero-sum condition:
\begin{equation}\label{app:local:pointwise-residual}
 \sum_b p_b\frac{[\Ima(\overline S_{ba}z_b^*Kh)]^2}{|z_b^*Kz_a|}
       -\lambda h^*Kh\le-c_{\rm intr}\|h\|^2,\qquad
 \Ima(z_a^*h)=0.
\end{equation}
Every columnwise $\xi_j$ is eligible since $z_j^*\xi_j=0$. Therefore,
for any certified upper $\lambda_+$,
\begin{equation}\label{app:local:spectral-row}
 H_{\rm spec}=\lambda_+(1+\kappa)-c_{\rm intr},\qquad
 H\leftarrow\min(H,H_{\rm spec}),\quad
 H_{\mathrm c}\leftarrow\min(H_{\mathrm c},H_{\rm spec}).
\end{equation}
The certificate checks $H_{\rm spec}\ge0$. Every summand in the row
is nonnegative, which justifies the cross-only bound. The coefficient
$M$ is not changed: a row bound does not imply a pointwise pair bound.
Nor is a mean curvature coefficient substituted for $c_{\rm intr}$.
Using \eqref{app:local:pointwise-residual} both here and in quadratic
descent is valid because both are pointwise consequences of the same
inequality; neither spends a separate portion of a deficit.

Hermitian symmetry gives $b_{ij}=-a_{ji}$, $c_{ji}=-c_{ij}$.
The full ordered radial term in \eqref{app:local:exact-pair} is consequently
\[
 \sum_{i,j}(t_i-t_j)a_{ij}^2/(r_js_{ij}).
\]
Put $W_i=\sum_j a_{ij}^2/s_{ij}$. Then $0\le W_i\le MX^2$,
$\sum_iW_i\le NHX^2$, hence $\sum_iW_i^2\le NMHX^4$.
Cauchy--Schwarz, followed by \eqref{app:local:shared-budget}, bounds this
term divided by $Nd^2$ by
\begin{equation}\label{app:local:radial-row}
 B_{\rm rad}\vartheta\sqrt{1-\vartheta^2},\qquad
 B_{\rm rad}=\frac{HT+\sqrt{MH}\,q_t}{\rho r_-}.
\end{equation}
Indeed the two sums with $|t_i|$ and $|t_j|$ are at most
$\sqrt{NMH}\,T_rX^2$ and $NHT_rTX$, respectively, before division by
$\rho$.

The full ordered linear--quadratic term is
$2\sum_{i,j}a_{ij}c_{ij}/(r_js_{ij})$. If
$E_b=\sum_{a(j)=b}e_j^2$, its cross part, by weighted Cauchy--Schwarz,
is at most
\[
 \frac{2g_{\mathrm c}}\rho\sqrt{NH_{\mathrm c}/h_{\mathrm c}}
 \sum_b E_b\sqrt{X^2-E_b}
 \le\frac{\sqrt3g_{\mathrm c}}\rho
        \sqrt{NH_{\mathrm c}/h_{\mathrm c}}X^3.
\]
The last inequality follows from
$\sum E_b^2\ge X^4/4$. The same-cluster part is at most
$2\sqrt{Np_+}m_{\mathrm s}g_{\mathrm s}X^3/(\rho h_{\mathrm s})$.
Set
\begin{gather}
 L_{\rm row}=\frac{\sqrt3}2g_{\mathrm c}\sqrt{H_{\mathrm c}/h_{\mathrm c}}
        +\sqrt{p_+}m_{\mathrm s}g_{\mathrm s}/h_{\mathrm s},\qquad
 J_{\max}=\max_t g_t^2/(2h_t),\nonumber\\
 A_{\rm mixed}=2L_{\rm mixed}q_t/\rho+J_{\max}q_t^2/\rho^2,
 \label{app:local:mixed-row}
\end{gather}
where $L_{\rm mixed}$ can be $L_{\rm row}$ or its minimum with the
independent weighted type bound described next. The second term in
$A_{\rm mixed}$ retains the full quadratic--quadratic term $c_{ij}^2$.
The same-cluster term in $L_{\rm row}$ is also retained.

\subsection{Complete scalar bounds and four-cluster summation}
\label{app:local:scalar-tail}

The elementary modulus identities below retain all orders of the
remainder; no assertion that its argument has modulus less than its
base is required. For $h\ne0$ and $s>0$ put
\[
 z=|h|/s,\qquad t=(1-|1+h/s|)/z,\qquad A_t=1-t^2.
\]
Reverse triangle gives $-1\le t\le1$. Squaring the defining identity
gives
\[
 \frac{\Rea h}{s}=-zt-z^2A_t/2,\qquad
 \frac{(\Ima h)^2}{s^2}=z^2A_t-z^3tA_t-z^4A_t^2/4.
\]
Consequently
\begin{equation}\label{app:local:scalar-identity}
 \frac{R_s(h)}s=z^3 f(t)+z^4 g(t),\qquad
 f(t)=\frac{t(1-t^2)}2,\quad g(t)=\frac{(1-t^2)^2}8.
\end{equation}
Replacing $t$ by $|t|$ only increases the right side. On $[0,1]$ put
$a_0=1/\sqrt3$. Then
\[
 f(a_0)-f(t)=\frac{(t-a_0)^2(t+2a_0)}2
 \ge a_0(t-a_0)^2,\qquad
 g(a_0)=1/18,\qquad |g'(t)|\le1/(3\sqrt3).
\]
The mean-value inequality followed by completing a square proves
\begin{equation}\label{app:local:unsigned-scalar}
 R_s(h)\le \frac{|h|^3}{3\sqrt3s^2}
       +\left(\frac1{18}+\frac Z{36\sqrt3}\right)\frac{|h|^4}{s^3},
 \qquad |h|/s\le Z.
\end{equation}
At $h=0$ the assertion is immediate.

Define the fixed mass coefficient $\zeta=(3/2)\sqrt{p_+}<1$.
For any nonnegative frozen pair of coefficients
$F_{\mathrm c}\ge F_{\mathrm s}$ set
\begin{equation}\label{app:local:mass-weights}
 F_{\rm eff}=(1-p_-)F_{\mathrm c}+p_-F_{\mathrm s},\quad
 F_{\rm mix}=\zeta F_{\mathrm c}+(1-\zeta)F_{\mathrm s},\quad
 F_{\rm pair}=(3F_{\mathrm c}+F_{\mathrm s})/4 .
\end{equation}
Here and below, ordering is checked for the chosen majorants, not
inferred from an ordering of the actual geometric quantities.
The ordered cubic cross sum satisfies
\[
 \sum_{a(i)\ne a(j)}e_i^2e_j\le\zeta\sqrt N X^3,\qquad
 \sum_{i,j}e_i^2e_j\le\sqrt N X^3.
\]
Indeed, if $x_a=\sqrt{\sum_{a(i)=a}e_i^2}/X$, then $\sum x_a^2=1$.
The first sum is at most
$\sqrt{Np_+}X^3(\sum x_a-\sum x_a^3)$.
The inequalities $\sum x_a\le2$ and
$\sum x_a^3\ge(\sum x_a)/4$ prove the claim. The latter follows from
$4\sum x_a^3\ge(\sum x_a^2)(\sum x_a)$; expanding gives a sum of
$(x_a-x_b)(x_a^2-x_b^2)\ge0$.
Also $\sum_{a(i)\ne a(j)}e_i^2e_j^2\le3X^4/4$, since
$\sum_aE_a^2\ge X^4/4$. Pure powers are weighted by the number of
cross partners, at most $N(1-p_-)$.
These proofs use arbitrary column energies and do not assume
constant perturbations in a cluster.

In particular the independent type bound used in
\eqref{app:local:mixed-row} is $(L_t)_{\rm mix}$ with
$L_t=m_tg_t/h_t$; take its minimum with $L_{\rm row}$ when
$L_{\mathrm c}\ge L_{\mathrm s}$.
For the scalar term of \eqref{app:local:exact-pair}, put
\begin{equation}\label{app:local:unsigned-types}
 v_t=\frac{g_tT}{2\rho},\quad D_t=m_t+v_t,\quad
 Z_t=\frac{2D_tT}{\rho h_t},\quad
 C_t=\frac{D_t^3}{3\sqrt3h_t^2},\quad
 B_t=\frac{D_t^4}{h_t^3}\left(\frac1{18}+\frac{Z_t}{36\sqrt3}\right).
\end{equation}
The argument of $R_s$ has modulus at most
$D_t(e_i/r_i+e_j/r_j)$, since
$e_ie_j/(r_ir_j)\le T(e_i/r_i+e_j/r_j)/(2\rho)$.
Keep $r_ir_j$ outside the scalar remainder until expanding the powers:
\begin{align*}
 r_ir_j(e_i/r_i+e_j/r_j)^3
 &=\frac{r_j}{r_i^2}e_i^3+\frac{r_i}{r_j^2}e_j^3
       +3e_i^2e_j/r_i+3e_ie_j^2/r_j,\\
 r_ir_j(e_i/r_i+e_j/r_j)^4
 &=\frac{r_j}{r_i^3}e_i^4+\frac{r_i}{r_j^3}e_j^4
       +4e_i^3e_j/r_i^2+4e_ie_j^3/r_j^2
       +6e_i^2e_j^2/(r_ir_j).
\end{align*}
Using \eqref{app:local:mass-weights} gives the full unsigned coefficients
\begin{align}
 A_3^{\,u}&=2RC_{\rm eff}T/\rho^2+6C_{\rm mix}q_t/\rho,\nonumber\\
 A_4^{\,u}&=2RB_{\rm eff}T^2/\rho^3
             +8B_{\rm mix}Tq_t/\rho^2+6B_{\rm pair}q_t^2/\rho^2,
 \label{app:local:unsigned-loss}\\
 A^{\,u}&=A_3^{\,u}+A_{\rm mixed}+A_4^{\,u},\qquad B^{\,u}=B_{\rm rad}.
 \nonumber
\end{align}
The necessary $C$ and $B$ orderings are part of the certificate guards.

\begin{lemma}[Shared-budget remainder]\label{app:local:loss}
Suppose one complete scalar comparison supplies coefficients $A,B\ge0$
such that the ordered remainder in \eqref{app:local:exact-pair}, divided
by $Nd^2$, is at most
$A\vartheta^2+B\vartheta\sqrt{1-\vartheta^2}$.
Then it is at most
\begin{equation}\label{app:local:loss-formula}
 {\rm Loss}(A,B)=\frac{A+\sqrt{A^2+B^2}}2.
\end{equation}
In particular \eqref{app:local:unsigned-loss} supplies such a comparison.
\end{lemma}
\begin{proof}
The maximum is the larger eigenvalue of
$\left(\begin{smallmatrix}A&B/2\\B/2&0\end{smallmatrix}\right)$.
Equations \eqref{app:local:radial-row},
\eqref{app:local:mixed-row} and the preceding power expansions establish
the unsigned assertion with the same $\vartheta$ in every term.
\end{proof}

\subsection{A signed remainder bound}
\label{app:local:signed}

Let $P_s(h)=-(\Rea h)(\Ima h)^2/(2s^2)$ be the scalar cubic.
Substitution in \eqref{app:local:scalar-identity} yields the exact identity
\begin{equation}\label{app:local:compact-identity}
 \frac{R_s(h)-P_s(h)}s=z^4H_z(t),\qquad
 H_z(t)=\frac{(1-t^2)(5t^2-1)}8
       +\frac{3z}{8}t(1-t^2)^2+\frac{z^2}{16}(1-t^2)^3 .
\end{equation}
This identity includes every order in $h$. Replacing $t$ by $|t|$
raises its right side, and $H_z(t)$ is nondecreasing in $z$ for
$0\le t\le1$. More generally choose a fixed sign weight
$\omega\in[0,1]$ and a rational $c_3\ge1/(3\sqrt3)$. For $Z>0$ let
\[
 W_{Z,\omega}(t)=
 \frac{(1-\omega)(f(t)-c_3)}Z+(1-\omega)g(t)+\omega H_Z(t).
\]
If $Q_\omega(Z)\ge\max_{0\le t\le1}W_{Z,\omega}(t)$, then
\begin{equation}\label{app:local:fractional-scalar}
 R_s(h)\le\omega P_s(h)+(1-\omega)c_3|h|^3/s^2
                       +Q_\omega(Z)|h|^4/s^3,\quad |h|/s\le Z.
\end{equation}
To verify it, subtract the first two terms, divide by $sz^4$, and use
\eqref{app:local:compact-identity}. The result is
$(1-\omega)(f(t)-c_3)/z+(1-\omega)g(t)+\omega H_z(t)$.
Since $f(t)\le c_3$, replacing $0<z\le Z$ by $Z$ increases it.
The zero case is direct. No condition $Z<1$ enters this proof.

The continuum maximum is bounded by rational Bernstein coefficients.
The ascending power coefficients of $H_Z$ are
\[
 \left(\frac{-2+Z^2}{16},\frac{3Z}{8},
 \frac{12-3Z^2}{16},-\frac{3Z}{4},
 \frac{-10+3Z^2}{16},\frac{3Z}{8},-\frac{Z^2}{16}\right).
\]
Add the displayed $f,g,c_3$ terms to obtain $W_{Z,\omega}$.
On a closed interval $[a,b]\subset[0,1]$, substitute
$t=a+(b-a)v$ and call the resulting power coefficients $c_j$.
Its degree-six Bernstein coefficients are
\[
 \beta_i=\sum_{j=0}^i c_j\frac{\binom ij}{\binom6j},\qquad 0\le i\le6.
\]
The identity
$W_{Z,\omega}(a+(b-a)v)=\sum_i\beta_i\binom6iv^i(1-v)^{6-i}$
and the nonnegative partition of unity bound the whole interval by
$\max_i\beta_i$. The certificate uses the $32$ closed equal intervals.
It checks $Q_\omega(Z)\ge0$ before substituting the nonnegative power
majorants in the ordered sums.

We next explain the pure-row cubic input; its sharpened version is
proved in the following subsection. For a fixed column $j$ choose a
unit vector $v_j\perp z_j$ and set
\[
 \ell_{ij}=\overline S_{ij}z_i^*Kv_j,\qquad
 q_{ij}=|\ell_{ij}|^2/s_{ij}^2 .
\]
Stationarity, with its full stress term, is
\begin{equation}\label{app:local:cubic-stationarity}
 \frac1N\sum_i\ell_{ij}=\lambda z_j^*Kv_j,\qquad
 |\lambda z_j^*Kv_j|\le\lambda_+\min\{\kappa r_+,m_{\mathrm s}\}.
\end{equation}
Its right side is generally nonzero. The Fourier identity
\[
 P_s(\ell u)=
 \Rea\{\ell^3u^3-|\ell|^2\ell|u|^2u\}/(8s^2)
\]
shows that a pure row bound follows from its first and third harmonics.
Given $q_{\mathrm c}\in[q_-,q_+]$, $q_{\mathrm s}\in[0,q_{\mathrm s,+}]$
and any fixed $c\ge0$, define
\begin{align}
 F_1(c)={}&c\lambda_+\min\{\kappa r_+,m_{\mathrm s}\}
 +(1-p_-)\max\{|q_--c|,|q_+-c|\}m_{\mathrm c}\nonumber\\
 &+p_+\max\{c,|q_{\mathrm s,+}-c|\}m_{\mathrm s},\nonumber\\
 F_3={}&(1-p_-)m_{\mathrm c}^3/h_{\mathrm c}^2
                  +p_+m_{\mathrm s}^3/h_{\mathrm s}^2,\qquad
 P_{\rm row}=(F_1(c)+F_3)/8 .\label{app:local:generic-cubic-row}
\end{align}
Subtract $c$ times \eqref{app:local:cubic-stationarity} from the first
harmonic and apply the triangle inequality; the third harmonic uses
its full typed triangle bound. Thus
\begin{equation}\label{app:local:pure-row}
 N^{-1}\sum_iP_{s_{ij}}(\alpha_{ij})\le P_{\rm row}e_j^3.
\end{equation}
Valid generic ratio intervals are
\begin{align*}
 q_-&=\max\left\{0,\frac{(1-\kappa)^2r_-^2}{(1/\sqrt3+d_g)^2}
                              -r_-^{-2}\right\},\\
 q_+&=\min\left\{\frac{m_{\mathrm c}^2}{h_{\mathrm c}^2},
       \frac{(1+\kappa)^2r_+^2}{h_{\mathrm c}^2}-r_+^{-2}\right\},
 \qquad q_{\mathrm s,+}=m_{\mathrm s}^2/h_{\mathrm s}^2 .
\end{align*}
They follow from the exact physical projection identity proved below.
The certificate can intersect these intervals with sharper valid
intervals and choose $c=2$ or their endpoint midpoint. Its chosen $c$
is fixed for the infinite-range argument.

To retain all mixed cubics, the symmetric trilinear form of $P_s$ obeys
$|\mathcal T_s(u,u,v)|\le|u|^2|v|/(3\sqrt3s^2)$.
For $u=x+iy$, the numerator norm is
$\sqrt{y^4+4x^2y^2}/6$, whose maximum at fixed $|u|$ gives this bound.
Expansion of the cubic, or integration of its gradient, gives
\[
 P_s(u+v+w)\le P_s(u)+P_s(v)+
 \frac{(|u|+|v|+|w|)^3-|u|^3-|v|^3}{3\sqrt3s^2}.
\]
All coefficients of the last polynomial are nonnegative.
For each pair type retain $v_t,D_t,Z_t$ from
\eqref{app:local:unsigned-types} and set
\begin{align}
 C_{{\rm base},t}&=c_3m_t^3/h_t^2,\nonumber\\
 C_{{\rm extra},t}&=
 c_3(3m_t^2v_t+3m_tv_t^2+v_t^3)/h_t^2
                      +(1-\omega)C_{{\rm base},t},\nonumber\\
 C_{{\rm full},t}&=c_3(m_t+v_t)^3/h_t^2,\qquad
 B_{{\rm signed},t}=D_t^4Q_\omega(Z_t)/h_t^3.
 \label{app:local:signed-types}
\end{align}
In particular $C_{\rm full}$ does not receive an additional
$(1-\omega)C_{\rm base}$ term. Outward rounding may increase $D_t,Z_t$;
the scalar certificate then uses that rounded radius.

The physical radial factors require a separate estimate. The two pure
ordered cubic terms combine to
$2\sum_jr_j^{-2}\sum_ir_iP_{s_{ij}}(\alpha_{ij})$.
One cannot pull a positive upper bound for $r_i$ out of this signed sum.
Write $r_i=1+t_i$ and apply \eqref{app:local:pure-row} only to the part
with coefficient one. With
\[
 C_{\rm rms}=\sqrt{(1-p_-)C_{{\rm base},{\mathrm c}}^2+
                         p_-C_{{\rm base},{\mathrm s}}^2},
 \qquad B_{\rm weight}=\frac{2\omega C_{\rm rms}Tq_d}{\rho^2r_-},
\]
Cauchy--Schwarz bounds the remaining part by
$B_{\rm weight}\vartheta^2\sqrt{1-\vartheta^2}$ after division by $Nd^2$.
Indeed $\sum_i|t_i|C_{{\rm base},t(i,j)}
\le\sqrt N C_{\rm rms}T_r$ and $\sum_j e_j^3\le TX^2$.
It is safe to replace $\vartheta^2$ by $\vartheta$ here.

Using \eqref{app:local:mass-weights}, the complete signed alternative is
\begin{align}
 A_3^{\,s}&=\frac{2\omega P_{\rm row}T}{\rho^2}
      +\frac{2R C_{{\rm extra},{\rm eff}}T}{\rho^2}
      +\frac{6C_{{\rm full},{\rm mix}}q_t}{\rho},\nonumber\\
 A_4^{\,s}&=\frac{2RB_{{\rm signed},{\rm eff}}T^2}{\rho^3}
      +\frac{8B_{{\rm signed},{\rm mix}}Tq_t}{\rho^2}
      +\frac{6B_{{\rm signed},{\rm pair}}q_t^2}{\rho^2},\nonumber\\
 A^{\,s}&=A_3^{\,s}+A_{\rm mixed}+A_4^{\,s},\qquad
 B^{\,s}=B_{\rm rad}+B_{\rm weight}.
 \label{app:local:signed-loss}
\end{align}
The certificate checks cross/same ordering for $C_{\rm base}$,
$C_{\rm extra}$, $C_{\rm full}$ and $B_{\rm signed}$, and nonnegative
quartic coefficients. Lemma~\ref{app:local:loss} then applies.
Taking the smaller of complete unsigned and complete signed bounds,
including the four fixed choices $\omega=1/4,1/2,3/4,1$, is valid.
Their cubic and scalar-tail pieces are not selected independently.

\subsection{Exact two-pair fibers and row cancellation}
\label{app:local:two-pair}

Here the first index of a fiber denotes its tangent center. Retain
$\mathcal T,d_a,\ell_a$ from \eqref{app:local:intrinsic-scalars}, and put
\begin{gather}
 t_{2a}=d_a^2+\sum_{b\ne a}\ell_a^2\ell_b^2p_ap_b,\qquad
 W_a=\frac{2t_{2a}}{p_a(\tr\mathcal T)^2},\qquad
 R_a=\frac{3\lambda-2+2p_a}{6p_a(2\lambda-1)}.
 \label{app:local:physical-fiber-data}
\end{gather}
Then $W_a=\|Kz_a\|^2$ and $R_a=\|z_a\|^2$.
Indeed $\mathcal T UU^*\mathcal T=\mathcal T^2$ since $UU^*$ projects
onto its range; its diagonal gives the formula for $W_a$.
The formula for $R_a$ is the candidate diagonal identity
$p_aR_a=1/4+(p_a-1/4)/[3(2\lambda-1)]$.
Complex dimension two implies that for every unit $v\perp z_a$,
\begin{equation}\label{app:local:exact-physical-fiber}
 m_{ab}^2:=|z_b^*Kv|^2=W_b-s_{ab}^2/R_a,\qquad
 q_{ab}=m_{ab}^2/s_{ab}^2 .
\end{equation}
This follows by subtracting the projection of $Kz_b$ on
$z_a/\|z_a\|$. In particular the modulus is independent of the choice
of unit vector in this one-dimensional complex fiber, without removing
its phase. For $a=b$ the overlap is $s_{aa}=d_a/(p_a\tr\mathcal T)$,
not the cross formula $\ell_a\ell_b/\tr\mathcal T$.

For the profile $C=1$ put
\begin{equation}\label{app:local:two-pair-parameter}
\begin{gathered}
 x=1/N,\quad t=\sqrt{3-8x^2},\quad
 a=\tfrac14-\tfrac x2,\quad b=\tfrac14+\tfrac x2,\\
 (p_1,p_2,p_3,p_4)=(a,a,b,b),\quad \lambda=(1+t)/4.
\end{gathered}
\end{equation}
Substitution in \eqref{alg:root-equation}, using $t^2+8x^2=3$,
proves the last identity by the uniqueness of the root in $(2/3,\alpha]$.
Reversing the signed count profile permutes the two pairs.
There are six ordered types: $\mathrm{aa},\mathrm{ab},\mathrm{ba},
\mathrm{bb},\mathrm{same}_a,\mathrm{same}_b$.
Here $\mathrm{aa}$ means two different clusters of mass $a$ and must
not be confused with $\mathrm{same}_a$.
The following explicit formulas make their interval verification
independent of a choice of spinor phases:
\begin{align}
 s_{\mathrm{aa}}&=\frac{t(2x+1)}{(2x-1)(8x^2-3)},&
 s_{\mathrm{ab}}&=\frac{-t}{8x^2-3},&
 s_{\mathrm{same}_a}&=
       -\frac{2tx+8x^2-3}{(2x-1)(8x^2-3)}. \label{app:local:sharp-s}
\end{align}
Set $V=4tx+20x^2+8x-3$ and
$J=-32tx^3-12tx^2+6tx+16x^4+56x^3-36x^2-24x+9$.
Then
\begin{align}
 m_{\mathrm{aa}}^2&=
 \frac{-2V}{(2x-1)(44x^2+12x-9)},\nonumber\\
 m_{\mathrm{ab}}^2&=
 \frac{2J}{(2x+1)(8x^2-3)(44x^2+12x-9)},\nonumber\\
 m_{\mathrm{same}_a}^2&=
 \frac{8x^2V}{(2x-1)(8x^2-3)(44x^2+12x-9)}, \label{app:local:sharp-m}\\
 q_{\mathrm{aa}}&=
 \frac{2(2x-1)(8x^2-3)V}{(2x+1)^2(44x^2+12x-9)},\nonumber\\
 q_{\mathrm{ab}}&=
 \frac{-2J}{(2x+1)(44x^2+12x-9)},\nonumber\\
 q_{\mathrm{same}_a}&=
 \frac{8x^2(-32tx^2+104x^3-36x^2-42x+9)}
 {9(2x-1)(2x+1)(44x^2+12x-9)}. \label{app:local:sharp-q}
\end{align}
All formulas for center $b$ are obtained by replacing $x$ by $-x$
and interchanging the labels $a,b$. These identities follow by
substitution in \eqref{app:local:physical-fiber-data} and
\eqref{app:local:exact-physical-fiber}, followed by multiplication by
the denominators and reduction using $t^2=3-8x^2$.
This supplies direct polynomial checks of each identity. The factor
$x^2$ in both same-type formulas is retained before interval evaluation.

For a center of mass $a$, denote the four complex forms in
\eqref{app:local:cubic-stationarity} by $\ell_0$ in its own cluster,
$\ell_1$ in the other $a$ cluster, and $\ell_2,\ell_3$ in the $b$
clusters. Exact stationarity reads
\[
 a\ell_0+a\ell_1+b\ell_2+b\ell_3=\lambda\ell_0.
\]
The two $b$ ratios are the same real number $q_{\mathrm{ab}}$ even
though their complex forms need not have the same phase. Thus the
first harmonic has the exact reduction
\begin{equation}\label{app:local:two-pair-harmonic}
 \sum_jp_jq_j\ell_j
 =a(q_{\mathrm{aa}}-q_{\mathrm{ab}})\ell_1+
 [a q_{\mathrm{same}_a}
                  +q_{\mathrm{ab}}(\lambda-a)]\ell_0.
\end{equation}
There is no assertion that the individual $b$ terms vanish.
Since $\lambda>a$,
\begin{align}
 F_{1a}&=a|q_{\mathrm{aa}}-q_{\mathrm{ab}}|m_{\mathrm{aa}}
    +[a q_{\mathrm{same}_a}+q_{\mathrm{ab}}(\lambda-a)]
                                          m_{\mathrm{same}_a},\nonumber\\
 F_{3a}&=a q_{\mathrm{same}_a}m_{\mathrm{same}_a}
         +a q_{\mathrm{aa}}m_{\mathrm{aa}}
         +2b q_{\mathrm{ab}}m_{\mathrm{ab}},\nonumber\\
 P_{\rm row}^{\,\rm pair}
   &=\max\{F_{1a}+F_{3a},F_{1b}+F_{3b}\}/8
 \label{app:local:two-pair-row}
\end{align}
is a valid bound in \eqref{app:local:pure-row}.
The third harmonic uses $|\ell_j|^3/s_j^2=q_jm_j$ exactly.
The $b$ center uses the exchanged formulas. In particular the
same-cluster stress term survives, and all physical residual vectors
are eligible. One may take the minimum with
\eqref{app:local:generic-cubic-row}.

For clarity, the finite continuum certificate for these formulas uses
the $512$ closed intervals
$[j/(180\cdot512),(j+1)/(180\cdot512)]$, $0\le j<512$.
On each interval $[u,v]$ use
$\sqrt{3-8v^2}\le t\le\sqrt{3-8u^2}$ with directed rational square roots.
Rational interval substitution in
\eqref{app:local:sharp-s}--\eqref{app:local:sharp-q}, for all six types,
checks that every denominator excludes zero. Geometric nonnegativity
permits clipping a lower bound for $m^2,q$ to zero; a negative upper
bound fails the check. Evaluate \eqref{app:local:two-pair-row} on each
same interval, retaining both centers.
The complete recorded bounds on $0\le x\le1/180$ include
\begin{align}
 m_{\mathrm c}^2&\le\frac{674866483913819}{10^{15}},&
 m_{\mathrm s}^2&\le\frac{1730019549}{62500000000000},\nonumber\\
 \frac{236851236625637}{125000000000000}
 &\le q_{\mathrm c}\le\frac{2109213874049103}{10^{15}},&
 q_{\mathrm s}&\le\frac{13969617457}{500000000000000},\nonumber\\
 P_{\rm row}^{\,\rm pair}&\le\frac{160219816186209}{10^{15}}.
 &&\label{app:local:sharp-uniform}
\end{align}
At the local endpoint $182$, every whole interval intersecting
$[0,1/182]$ is retained, including a crossing boundary interval.
The slightly sharper pure-row cap used there is
$32029714476747/200000000000000$.
This is an interval cover, not a maximum over sampled configurations.
Only the pure-row term is replaced in \eqref{app:local:signed-loss};
$B_{\rm weight}$ and every other remainder term are still present.

\subsection{Propagation and equality}
\label{app:local:closure}

We now extend the local comparison from a certified endpoint to all
later integers of the same profile. On each closed localization region,
fix its $U,W,L_U,L_P,\bar E$ and the simultaneous candidate
coefficients $B,C_{\mathrm{pt}}$ supplied by Sections~\ref{app:global} and
\ref{app:candidate}. At an accepted endpoint $n$, put $e_n=m_*/n$ and
$\sigma_n^2=C/n^2$. Let $\kappa_*$ be the outward rational upper enclosure
of $n$ times the explicit stress bound in \eqref{app:candidate:kappa}:
\[
 \kappa_*\ge
 2\sqrt C\sqrt{1/3+4e_n/b_0+4\sigma_n^2/b_0^2}
       +\frac{2ne_n^2}{1/15-e_n},\qquad b_0=109/100.
\]
This coefficient is distinct from the Bloch-Gram coefficient $k_0$ in
Section~\ref{app:candidate}. We retain that accepted enclosure throughout
propagation. At this endpoint, valid physical coefficients are
\begin{equation}\label{app:local:frozen-physical}
 s_n=\sqrt{13U/(4n)},\qquad
 A_R=\frac{13U}4+L_UW\frac{1+s_n}{(1-s_n)^2},\qquad
 P_R=\frac{13U}4+L_PW\frac{1+s_n}{(1-s_n)^2}.
\end{equation}
These follow from the radial/angular identity, with the relative
angular bound $L_Uw/[N(1-s)^2]$ and its pointwise analogue. Their minimum
with an independently certified bound in the original variables is
also valid. The partition and the closed count cover remain fixed, so
the count exclusions apply at every later local endpoint.

At an accepted endpoint, all of the following are checked by outward
rational inequalities: positive masses and overlap floors; the
candidate point domains; every fine phase guard and the separate
closest-gauge rotation guard; the mean-coordinate phase, projection,
and positivity guards; the intrinsic residual lower bound; positive
radial multipliers; both finite-normal branches; the nonnegative
quartic and cross/same ordering guards; and the full remainder margin.
The square root of a positive rational is enclosed by rational numbers
whose squares straddle it, so these checks require no floating-point
eigenvalue or optimizer. Section~\ref{app:certificates} specifies
the interval records, arithmetic and source binding.

\begin{lemma}[Frozen endpoint propagation]\label{app:local:propagation}
Suppose all stated guards and a strict complete local margin hold on
one closed region at endpoint $n$, with intrinsic and exact-fiber
enclosures valid on $0\le1/N\le1/n$. Then the same margin lower bound
holds on that region for every later admissible member of its profile.
\end{lemma}
\begin{proof}
For $N\ge n$, retain the fixed profile deviations $\delta_a$ and use
the actual masses $p_a=1/4+\delta_a/N$. The candidate is reconstructed
at those masses; it is not obtained by scaling a finite-endpoint vector
configuration. Its errors obey
\[
 e=m_*/N,\qquad \eta\le\sqrt{C_{\mathrm{pt}}C}/N,\qquad
 \kappa\le\kappa_*/N,\qquad \tau\le(2/3)C/N^2 .
\]
For the fixed count profile, the stress estimate follows because $N$
times the explicit bound in \eqref{app:candidate:kappa} is nondecreasing in $e=m_*/N$ and
$\sigma^2=C/N^2$; its value for $N\ge n$ is therefore bounded by the
accepted enclosure $\kappa_*$.
Freeze every favorable lower coefficient and every unfavorable upper
coefficient of the accepted endpoint. Here is the justification for
each part of this operation.

First, the global exclusions remain strict and the regional fine
bounds remain valid by Section~\ref{app:global}. In
\eqref{app:local:frozen-physical}, $s_n$ and $(1+s_n)/(1-s_n)^2$
decrease as the endpoint increases. The candidate errors decrease,
$r_-$ increases, and $r_+$ decreases. Thus all angular transport
majorants, $T,q_t,q_d$ and $R$ have their endpoint upper bounds;
$\rho$ has its endpoint lower bound. Both actual-rotation guards remain
valid. No difference of decreasing functions is asserted to be
monotone.

Second, each nonnegative perturbation, phase-defect and projection
majorant in the mass and stress mean estimates decreases when these
candidate errors decrease. Their favorable metric factors increase.
Hence the accepted physical mean lower bound remains valid. The
intrinsic $c_K$ was certified on the entire reciprocal-$N$ domain,
and multiplication by $1-\kappa$ preserves its endpoint lower bound.
The original physical normal estimates similarly retain their fixed
upper $J,L$. The normal radius $d_c/(\sqrt2N)$ shrinks. Therefore
the supremum in \eqref{app:local:normal-optimization} cannot increase:
its coefficient bounds are frozen and its domain only shrinks.
This argument does not keep a particular active maximizing branch.

Third, all generic projected and phase-row majorants have the displayed
endpoint bounds. The exact two-pair inequalities cover the whole
reciprocal domain, without needing monotonicity of an individual
ratio. Freeze $\lambda_+$ and $c_K$ separately in the spectral row
bound; then
$\lambda_+(1+\kappa)-(1-\kappa)c_K$ has its endpoint upper bound.
The cross-only row shares this bound, while the independent pair
coefficient $M$ is retained.

Finally freeze every scalar radius cap $Z_t$, sign weight $\omega$,
Bernstein upper bound $Q_\omega(Z_t)$, harmonic center $c$, and the
separate positive cross/same type coefficients. All later increments
lie in the same certified disks. The nonnegative type polynomials and
the signed physical-weight estimate give the endpoint upper
coefficients in \eqref{app:local:signed-loss}; the unsigned bound
propagates in the same way. The ordering checks concern these fixed
majorants. No monotonicity of a rounded Bernstein output or of a
selected numerical branch is needed.
\end{proof}

The rational certificates in Section~\ref{app:certificates:acceptance}
cover all $7320$ budget cells by strict count exclusions and $228$
accepted closed local regions. On every accepted region, the certified
inequality \eqref{app:certificates:accepted-margin} gives
$g_{\rm normal}-{\rm Loss}>1/10000$.
The detailed coverage totals and rational margins are recorded there.

\begin{proposition}[Uniform strict local comparison]
\label{app:local:strict}
For every nondivisible integer $N\ge180$ and every localized competitor
with its actual closest candidate and gauge,
\begin{equation}\label{app:local:strict-inequality}
 \Gamma_K(Y)\le N^2\lambda-\frac{N}{10000}\|Y-Z\|_F^2 .
\end{equation}
\end{proposition}
\begin{proof}
If $d=0$, stationarity gives equality in
\eqref{app:local:strict-inequality} directly. Otherwise, since
$D\mathcal L(Z)=0$, the second-order Taylor polynomial of
$\mathcal L(Z+\Delta)-\mathcal L(Z)$ is $Q(\Delta)$.
The quadratic estimate is at most $-Ng_{\rm normal}d^2$, and the
ordered Taylor remainder is at most $N\,{\rm Loss}\,d^2$.
Every retained closed region has
$g_{\rm normal}-{\rm Loss}>1/10000$ by the rational checks just stated.
The preceding propagation argument proves the entire subsequent
range. Odd integers begin at $181$; residue-two integers begin at
$182$, so the local endpoint $182$ for $C=1$ leaves no integer gap.
\end{proof}

The comparison applies to every localized kernel, with no restriction
on the number of distinct column directions.
In the ordered global trace chain, equality in the optimal bound
therefore forces $d=0$ by Proposition~\ref{app:local:strict}.
Strict integer balancing then forces the balanced profile, and every
nonzero candidate stress overlap fixes its corresponding matrix entry.
This gives the unitary, individual-phase, permutation and conjugation
classification in the main theorem. The divisible
integers, including $180$, are handled by the exact divisible case of
the global theorem. Hence $180$ is a sufficient threshold.
